% ============================================================
\chapter{Stability Theorems}
\label{ch:stability}
% ============================================================

Persistent homology would be a mere mathematical curiosity if it were not \emph{stable}: a small change in the data must produce a small change in the persistence diagram. This fundamental result, proved by David Cohen-Steiner, Herbert Edelsbrunner, and John Harer in 2007, is the pillar on which all applied TDA rests. It states that the bottleneck distance between two persistence diagrams is bounded by the $L^\infty$ distance between the filtration functions. Without this theorem, the slightest noise in the data could completely overturn the topological descriptors, rendering them useless in practice.

\begin{intuition}
Stability theorems are the theoretical foundation of TDA: they guarantee
that small perturbations of the data produce small changes in
persistence diagrams. Without stability, topological descriptors would
be useless in practice.
\end{intuition}

% ============================================================
\section{Distances between Diagrams: Review}
% ============================================================

Before stating the stability theorems, we recall the distances used.
These will be developed in Chapter~\ref{ch:distances}.

\begin{definition}[Bottleneck Distance]
Let $D_1, D_2$ be two persistence diagrams. The \emph{bottleneck
distance} is:
\[
  d_B(D_1, D_2) = \inf_{\gamma} \sup_{x \in D_1}
  \norm{x - \gamma(x)}_\infty
\]
where the infimum is over all bijections $\gamma : D_1 \to D_2$
(recall that diagrams include the diagonal with infinite multiplicity).
\end{definition}

\begin{definition}[Wasserstein Distance]
The \emph{$q$-Wasserstein distance} ($1 \leq q < \infty$) is:
\[
  W_q(D_1, D_2) = \left(\inf_{\gamma} \sum_{x \in D_1}
  \norm{x - \gamma(x)}_\infty^q \right)^{1/q}.
\]
\end{definition}

\begin{keyformulas}[title=\textbf{Relationship between Distances}]
\[
  d_B(D_1, D_2) = W_\infty(D_1, D_2) \leq W_q(D_1, D_2) \leq
  W_p(D_1, D_2) \quad \text{for } p \leq q.
\]
\end{keyformulas}

% ============================================================
\section{Bottleneck Stability Theorem}
% ============================================================

The fundamental theorem of TDA is due to Cohen-Steiner, Edelsbrunner,
and Harer (2007).

\begin{definition}[Tame Function]
A function $f : X \to \R$ is called \emph{tame} if the sublevel sets
$f^{-1}((-\infty, r])$ have finite-dimensional homology for all $r$,
and the number of critical values is finite.
\end{definition}

\begin{theorem}[Bottleneck Stability --- Cohen-Steiner, Edelsbrunner, Harer, 2007]
Let $f, g : X \to \R$ be two tame functions on a triangulable
topological space $X$. Then:
\[
  d_B(\mathrm{Dgm}(f), \mathrm{Dgm}(g))
  \leq \norm{f - g}_\infty.
\]
\end{theorem}

\begin{proof}[Proof sketch]
The proof relies on three ingredients:
\begin{enumerate}
  \item \textbf{Interpolation}: consider the family
        $f_t = (1-t)f + tg$ for $t \in [0, 1]$.
  \item \textbf{Box lemma}: if the diagrams of $f$ and $g$ differ in a
        region, then $\norm{f - g}_\infty$ is at least as large as the
        size of that region.
  \item \textbf{Matching construction}: build a bijection $\gamma$
        between diagrams by tracking points through the interpolation.
\end{enumerate}
See Cohen-Steiner, Edelsbrunner, Harer (2007) for the complete proof.
\end{proof}

\begin{remark}
This theorem is tight: there exist functions for which equality is
achieved.
\end{remark}

% ============================================================
\section{Stability for Point Clouds}
% ============================================================

\begin{corollary}[Vietoris--Rips Stability]
Let $X, Y \subset \R^d$ be two finite point clouds. Let $d_H(X, Y)$
be the Hausdorff distance. Then:
\[
  d_B(\mathrm{Dgm}(\mathrm{VR}(X)), \mathrm{Dgm}(\mathrm{VR}(Y)))
  \leq 2\, d_H(X, Y).
\]
\end{corollary}

\begin{definition}[Hausdorff Distance]
The \emph{Hausdorff distance} between two compact sets $X, Y$ in a
metric space $(M, d)$ is:
\[
  d_H(X, Y) = \max\left(
    \sup_{x \in X} \inf_{y \in Y} d(x, y),\;
    \sup_{y \in Y} \inf_{x \in X} d(x, y)
  \right).
\]
\end{definition}

\begin{proof}
Let $f_X(z) = \inf_{x \in X} d(z, x)$ be the distance function to $X$,
and similarly for $f_Y$. Then
$\norm{f_X - f_Y}_\infty = d_H(X, Y)$.
The persistence diagrams of the sublevel set filtrations of $f_X$ and
$f_Y$ coincide (up to constants) with those of the Rips filtrations.
The result follows from the bottleneck stability theorem, with the
factor~$2$ arising from the parametrization convention.
\end{proof}

% ============================================================
\section{Wasserstein Stability}
% ============================================================

\begin{theorem}[Wasserstein Stability --- Chazal et al., 2009]
Let $f, g : X \to \R$ be two tame functions. For all $q \geq 1$ and
$p \in \N$, if the diagrams $\mathrm{Dgm}_p(f)$ and
$\mathrm{Dgm}_p(g)$ have at most $N$ off-diagonal points, then:
\[
  W_q(\mathrm{Dgm}_p(f), \mathrm{Dgm}_p(g))
  \leq C(N, q) \cdot \norm{f - g}_\infty
\]
where $C(N, q) = N^{1/q}$ for the simplest bound.
\end{theorem}

\begin{remark}
Unlike bottleneck stability, Wasserstein stability depends on the
number of points in the diagram. Sharper bounds exist under additional
hypotheses.
\end{remark}

% ============================================================
\section{Stability of Representations}
% ============================================================

\begin{theorem}[Stability of Persistence Landscapes]
Persistence landscapes are $1$-Lipschitz with respect to the
bottleneck distance:
\[
  \norm{\lambda_k^{(1)} - \lambda_k^{(2)}}_\infty
  \leq d_B(D_1, D_2)
\]
for all $k \geq 1$.
\end{theorem}

\begin{proposition}[Stability of Persistence Images]
Persistence images are Lipschitz with respect to the $1$-Wasserstein
distance: there exists $C > 0$ (depending on $\sigma$ and the weight
function) such that:
\[
  \norm{\mathrm{PI}(D_1) - \mathrm{PI}(D_2)}_F
  \leq C \cdot W_1(D_1, D_2).
\]
\end{proposition}

% ============================================================
\section{Generalized Stability Theorem}
% ============================================================

\begin{definition}[Interleaving Distance]
Two persistence modules $\mathbf{U}$ and $\mathbf{V}$ are
\emph{$\varepsilon$-interleaved} if there exist natural morphisms
$\varphi : \mathbf{U} \to \mathbf{V}[\varepsilon]$ and
$\psi : \mathbf{V} \to \mathbf{U}[\varepsilon]$ such that the
composites $\psi[\varepsilon] \circ \varphi$ and
$\varphi[\varepsilon] \circ \psi$ coincide with the
$2\varepsilon$-shift structure maps.

The \emph{interleaving distance} is:
\[
  d_I(\mathbf{U}, \mathbf{V}) = \inf\{\varepsilon \geq 0 :
  \mathbf{U} \text{ and } \mathbf{V} \text{ are }
  \varepsilon\text{-interleaved}\}.
\]
\end{definition}

\begin{theorem}[Algebraic Stability Isometry]
For any p.f.d. persistence modules:
\[
  d_B(\mathrm{Dgm}(\mathbf{U}), \mathrm{Dgm}(\mathbf{V}))
  = d_I(\mathbf{U}, \mathbf{V}).
\]
\end{theorem}

\begin{remark}
This theorem, due to Chazal, Cohen-Steiner, Glisse, Guibas, and Oudot
(2009) and Bauer--Lesnick (2015), shows that the bottleneck distance is
the ``right'' distance on persistence diagrams: it coincides exactly
with the natural notion of proximity for persistence modules.
\end{remark}

% ============================================================
\section{Applications of Stability}
% ============================================================

\subsection{Topological Inference}

\begin{theorem}[Topological Inference --- Niyogi, Smale, Weinberger, 2008]
Let $\mathcal{M}$ be a compact submanifold of $\R^d$ with reach
$\tau > 0$. If $X$ is an $\varepsilon$-dense sample of $\mathcal{M}$
with $\varepsilon < \tau / 2$, then the \v{C}ech filtration of $X$ at
parameter $r \in (\varepsilon, \tau - \varepsilon)$ has the same
homotopy type as $\mathcal{M}$.
\end{theorem}

\subsection{Robustness to Noise}

\begin{proposition}
If $X$ is a point cloud sampled from a manifold $\mathcal{M}$ and $X'$
is a noisy version with $d_H(X, X') \leq \delta$, then:
\[
  d_B(\mathrm{Dgm}(X), \mathrm{Dgm}(X')) \leq 2\delta.
\]
Topological features with persistence $> 4\delta$ in the diagram of
$X'$ correspond to genuine features of $\mathcal{M}$.
\end{proposition}

% ============================================================
\section{Numerical Verification}
% ============================================================

\begin{codeblock}[title=\textbf{Numerical Verification of Stability}]
\begin{minted}{python}
import numpy as np
from ripser import ripser
from persim import bottleneck

# Exact circle
theta = np.linspace(0, 2*np.pi, 100, endpoint=False)
X = np.column_stack([np.cos(theta), np.sin(theta)])

# Noisy versions with increasing epsilon
epsilons = [0.01, 0.05, 0.1, 0.2, 0.5]
dgm_X = ripser(X, maxdim=1)['dgms']

for eps in epsilons:
    noise = eps * np.random.randn(*X.shape)
    X_noisy = X + noise
    dgm_Y = ripser(X_noisy, maxdim=1)['dgms']

    # Bottleneck distance in H_1
    d_bn = bottleneck(dgm_X[1], dgm_Y[1])

    # Hausdorff distance (approximation)
    from scipy.spatial.distance import directed_hausdorff
    d_h = max(directed_hausdorff(X, X_noisy)[0],
              directed_hausdorff(X_noisy, X)[0])

    print(f"eps={eps:.2f}: d_B={d_bn:.4f}, "
          f"d_H={d_h:.4f}, "
          f"ratio={d_bn/d_h:.4f}")
\end{minted}
\end{codeblock}

% ============================================================
\section{Limitations of Stability}
% ============================================================

\begin{warning}[title=\textbf{Limitations}]
\begin{itemize}
  \item Bottleneck stability is an \emph{upper} bound: the diagram may
        change much less than the bound suggests.
  \item The Wasserstein bound depends on the number of points: for
        diagrams with many low-persistence points, the bound can be
        loose.
  \item Stability concerns the \emph{diagrams}, not necessarily the
        \emph{vectorized representations} (although these often inherit
        stability properties).
\end{itemize}
\end{warning}

% ============================================================
\section{Exercises}
% ============================================================

\begin{exercise}
Show that the bottleneck distance satisfies the triangle inequality.
\end{exercise}

\begin{exercise}
Construct an example of two functions $f, g : [0, 1] \to \R$ such that
$d_B(\mathrm{Dgm}(f), \mathrm{Dgm}(g)) = \norm{f - g}_\infty$ (show
that the bound is tight).
\end{exercise}

\begin{exercise}
Numerically verify bottleneck stability by generating point clouds on a
torus with different noise levels. Plot $d_B$ versus $d_H$ and verify
linearity.
\end{exercise}

\begin{exercise}
Prove that if $\mathbf{U}$ and $\mathbf{V}$ are
$\varepsilon$-interleaved, then
$d_B(\mathrm{Dgm}(\mathbf{U}), \mathrm{Dgm}(\mathbf{V}))
\leq \varepsilon$.
\end{exercise}

\begin{exercise}
Let $X$ be a cloud of $n$ i.i.d. points on a compact manifold
$\mathcal{M}$. Show that $d_H(X, \mathcal{M}) \to 0$ almost surely as
$n \to \infty$, and deduce the convergence of persistence diagrams.
\end{exercise}
